%LTeX: language=fr-FR
\documentclass[12pt,a4paper]{article}
\usepackage[utf8]{inputenc}
\usepackage[T1]{fontenc}
\usepackage[french]{babel}
\usepackage[hmargin=1.7cm, vmargin=1.3cm]{geometry}
\usepackage{amsmath,amssymb}
\usepackage{graphicx}
\usepackage{enumitem}
\usepackage{tcolorbox}
\usepackage{xcolor}
\usepackage{multicol}
\usepackage{makecell}
\usepackage{tasks}

\tcbuselibrary{skins}

\definecolor{bandeau}{HTML}{E8EEF7}
\definecolor{titrebleu}{HTML}{1A4FA0}

% Bandeau de titre pour chaque exercice
\newtcbox{\exobox}{
  colback=bandeau, colframe=bandeau, boxrule=0pt,
  arc=0pt, left=4pt, right=4pt, top=3pt, bottom=3pt,
  fontupper=\bfseries, width=\linewidth, halign=left
}

% \exercice{numéro}{titre}{barème}
\newcommand{\exercice}[3]{%
  \vspace{0.6em plus 1fill}
  \noindent\exobox{Exercice #1 --- #2 \hfill (#3 points)}\par
  \vspace{0.4em}
}

\pagestyle{empty}
\setlist{topsep=3pt, itemsep=4pt, parsep=0pt}
\setlength{\parskip}{3pt}
\setlength{\parindent}{0pt}

\begin{document}

\begin{center}
  {\large\bfseries Contrôle de rattrapage --- Tableaux croisés et probabilités conditionnelles}\\[0.4em]
  {\itshape 1\textsuperscript{ère} ST2S \quad $\bullet$ \quad Durée~: 55 minutes \quad $\bullet$ \quad Calculatrice autorisée}
\end{center}

\vspace{0.3em}
\noindent\textit{La qualité de la rédaction et la justification des réponses seront prises en compte dans la notation.}

% ================= Partie A =================
\exercice{1}{Automatismes}{6}
Les trois questions sont indépendantes. Détailler les calculs.
\begin{enumerate}
  \item Le prix d'une place de cinéma est passé de $8$~euros à $10$~euros. Il a augmenté de~:

  \textbf{A.}~$2\,\%$ \hfill \textbf{B.}~$25\,\%$ \hfill \textbf{C.}~$20\,\%$ \hfill \textbf{D.}~$80\,\%$ \hspace*{2cm}
  \item Déterminer la solution de l'équation $-6x+4=-17$.
  \item Calculer $A = \dfrac{3}{1-\dfrac{1}{4}}$. Donner le résultat sous forme d'un entier ou d'une fraction irréductible.
\end{enumerate}

% ================= Partie B =================
\exercice{2}{Test de dépistage}{9}
Un laboratoire propose un test de dépistage rapide d'une maladie. On a testé $200$ patients, parmi lesquels $50$ sont malades. Parmi les patients malades, $80\,\%$ ont un test positif. Au total, $50$ patients ont un test positif.
\begin{enumerate}
  \item Recopier et compléter le tableau ci-dessous~:
  \begin{center}
  \renewcommand{\arraystretch}{1.2}
  \begin{tabular}{|c|c|c|c|}
  \hline
   & Patient malade & Patient non malade & Total \\ \hline
  Test positif & & & 50 \\ \hline
  Test négatif & & & \\ \hline
  Total & 50 & & 200 \\ \hline
  \end{tabular}
  \end{center}
\end{enumerate}
On choisit un patient au hasard parmi les $200$. On note $M$ l'événement \og le patient est malade\fg{} et $T$ l'événement \og le test du patient est positif\fg{}.
\begin{enumerate}[resume]
  \item Calculer les probabilités suivantes~:
  \begin{tasks}[label=\alph*), label-width=1.5em](4)
    \task $P(M)$
    \task $P(T)$
    \task $P_M(T)$
    \task $P(M\cap\overline{T})$
  \end{tasks}
  \item Le test du patient est positif. Quelle est la probabilité qu'il soit malade~?
  \item Décrire par une phrase l'événement $M\cup T$, puis calculer sa probabilité.
  \item Le test se trompe lorsqu'il est positif pour un patient non malade, ou négatif pour un patient malade. \\
  \textbf{Vrai ou faux~?} \og La probabilité que le test se trompe est inférieure à $15\,\%$.\fg{} Justifier.
\end{enumerate}

\exercice{3}{Construire un tableau croisé}{5}
Un hôpital interroge 100 agents sur leur horaire de travail et leur vaccination contre la grippe. Les résultats sont les suivants~:
\begin{multicols}{2}
\begin{itemize}[label=$\bullet$]
  \item 40 agents de jour sont vaccinés~;
  \item 10 agents de nuit sont vaccinés~;
  \item 20 agents de nuit ne sont pas vaccinés~;
  \item 30 agents de jour ne sont pas vaccinés.
\end{itemize}
\end{multicols}
\begin{enumerate}
  \item Construire un tableau croisé d'effectifs correspondant à la situation. \\
  \textit{Aide~: on mettra en lignes l'horaire de travail (jour ou nuit) et en colonnes la vaccination.}
  \item Quelle est la proportion d'agents non vaccinés~?
  \item Calculer le pourcentage d'agents de nuit parmi les agents vaccinés.
\end{enumerate}

\end{document}
